% archon:covers AlgebraicJacobian/RiemannRoch/WeilDivisor.lean
\chapter{Weil divisors on a smooth proper curve (RR.1)}
\label{chap:RiemannRoch_WeilDivisor}
% SOURCE: Hartshorne, \emph{Algebraic Geometry}, II \S6, pp.~130--137 +
% IV.1, pp.~294--296 (read from references/hartshorne-algebraic-geometry.pdf).
% Stacks pointers (tag-level, not verbatim-quoted): 02RW (divisors),
% 02ME (order at a point on a regular 1-dimensional scheme),
% 0BE0 (degree of a divisor on a curve), 0BE3 (principal divisors have
% degree zero on a complete nonsingular curve).
\section{Setup and motivation}
This chapter is the first of four sub-build chapters \texttt{RR.1}--\texttt{RR.4}
that close the project's headline \emph{RR bridge}
\[
  \mathtt{genusZero\_curve\_iso\_P1}
  \,\colon\,
  \text{a smooth proper geometrically irreducible curve \(C / \bar k\) with
  \(g(C) = 0\) is isomorphic to \(\mathbb P^1_{\bar k}\).}
\]
The bridge as proved in Hartshorne (Example~IV.1.3.5) chains together:
divisor of a closed point on a curve (\texttt{RR.1}, this chapter); the
Riemann--Roch dimension formula \(\ell(D) = \deg D + 1\) for \(\deg D \geq 0\) on
a genus-\(0\) curve (\texttt{RR.2}); the global-sections argument identifying
\(\dim_k H^0(C, \mathcal O_C(P)) = 2\) via \(1\) and a non-constant function
(\texttt{RR.3}); the ``rational curve \(\Rightarrow \cong \mathbb P^1\)''
classification recovered as a degree-\(1\) morphism into \(\mathbb P^1\) via
\texttt{Proj.fromOfGlobalSections} (\texttt{RR.4}).
Mathlib ships no \texttt{WeilDivisor} on a scheme. Adjacent pieces exist, but
none of them is the data structure we need: \texttt{MeromorphicOn.divisor}
(\texttt{Mathlib.Analysis.Meromorphic.Divisor}) is the order function of a
meromorphic function valued in a \emph{normed} field, not a formal sum of
codimension-\(1\) cycles on a scheme; \texttt{CommRing.Pic}
(\texttt{Mathlib.RingTheory.PicardGroup}) is the ring-level Picard group of
invertible modules under tensor product, with no scheme-level analogue;
\texttt{AlgebraicGeometry.Scheme.RationalMap}
(\texttt{Mathlib.AlgebraicGeometry.Birational.RationalMap}) packages rational
maps but not divisors. Consequently this chapter, and the three siblings, are
\emph{project-bespoke}: the only Mathlib tools they rely on are
discrete-valuation rings, the function field, and the ambient scheme API.
This chapter establishes the formal-sum data type
\(\mathrm{Div}(X) = \bigoplus_{\substack{Z \subset X \\ \mathrm{codim} 1}} \mathbb Z\) on a
Noetherian integral scheme \(X\) satisfying Hartshorne's condition \((*)\), the
principal-divisor homomorphism \(\mathrm{div} \colon K(X)^{\times} \to \mathrm{Div}(X)\) on a
curve, the degree map \(\deg \colon \mathrm{Div}(C) \to \mathbb Z\) on a smooth proper
curve over an algebraically closed field, and the linear-equivalence
relation. The \emph{degree-zero of a principal divisor} statement
(\ref{thm:principal_deg_zero}) is included for completeness; its proof
decomposes into a separate sub-build inside \texttt{RR.1} (it uses the
auxiliary lemma that any non-constant rational function on \(C\) defines a
finite morphism \(C \to \mathbb P^1\)). The line-bundle construction $\mathcal
O_C(D)$ and the associated cohomology computation $H^0(C, \mathcal O_C(D))$
live in the sibling chapter \texttt{RR.3}; the Riemann--Roch dimension
formula lives in \texttt{RR.2}; the genus-\(0\) classification proper lives in
\texttt{RR.4}.
Throughout this chapter, the standing convention is Hartshorne's condition
\((*)\):
% SOURCE QUOTE: "\((\ast)\) \(X\) is a noetherian integral separated scheme which
% is regular in codimension one."
\textit{Source: Hartshorne, II \S6, p.~130, the conditional \((*)\).}
\[
  (\ast)\qquad
  X \text{ is a noetherian integral separated scheme which is regular in
  codimension one.}
\]
For the degree map and for principal-divisor statements on a curve, we
additionally assume \(X\) is a complete (i.e.\ proper) nonsingular curve over
an algebraically closed field \(\bar k\), where Hartshorne's condition \((*)\) is
automatic.
\section{Codim-\(1\) cycle group}
\begin{definition}
\leanok
[Prime divisor of a scheme]
  \label{def:prime_divisor}
  \lean{AlgebraicGeometry.Scheme.PrimeDivisor}
  % SOURCE: [Hartshorne], II.6, p.~130 (definition of "prime divisor" under
  % condition \((*)\)) (read from references/hartshorne-algebraic-geometry.pdf,
  % PDF page 147). Stacks pointer: tag 02RW (divisors).
  % SOURCE QUOTE: "Definition. Let \(X\) satisfy \((\ast)\). A \emph{prime
  % divisor} on \(X\) is a closed integral subscheme \(Y\) of codimension one."
  \textit{Source: Hartshorne, II.6, p.~130 (definition of prime divisor).}
  Let \(X\) satisfy \((*)\). A \emph{prime divisor} on \(X\) is a closed integral
  subscheme \(Y \subset X\) of codimension one. Equivalently, by the standard
  generic-point correspondence (Hartshorne II \S6, p.~130, ``let $\eta \in
  Y$ be its generic point\ldots''), $Y$ is determined by its generic point
  \(\eta \in X\): closed integral subschemes of codimension one of \(X\) are in
  natural bijection with points \(x \in X\) such that the closure
  \(\overline{\{x\}} \subset X\) has codimension one. The closure
  \(\overline{\{x\}}\), equipped with the \emph{reduced induced subscheme
  structure}, is automatically integral (irreducible, because the closure
  of a single point is irreducible, and reduced by construction); the only
  non-automatic condition is the codimension.
\end{definition}
\begin{definition}
\leanok
[Codim-\(1\) cycle group / Weil divisor group of a scheme]
  \label{def:codim1_cycles}
  
  \lean{AlgebraicGeometry.Scheme.WeilDivisor}
  \uses{def:prime_divisor}
  % SOURCE: [Hartshorne], II.6, p.~130 (definition of "Weil divisor" under
  % condition \((*)\)) (read from references/hartshorne-algebraic-geometry.pdf,
  % PDF page 147). Stacks pointer: tag 02RW.
  % SOURCE QUOTE: "A \emph{Weil divisor} is an element of the free abelian
  % group \(\operatorname{Div} X\) generated by the prime divisors. We write a
  % divisor as \(D = \sum n_i Y_i\), where the \(Y_i\) are prime divisors, the
  % \(n_i\) are integers, and only finitely many \(n_i\) are different from
  % zero. If all the \(n_i \geq 0\), we say that \(D\) is \emph{effective}."
  \textit{Source: Hartshorne, II.6, p.~130 (Weil divisor group).}
  Let \(X\) be a Noetherian integral separated scheme that is regular in
  codimension one (i.e.\ \(X\) satisfies \((*)\)). The \emph{Weil divisor group}
  \[
    \mathrm{Div}(X) \;=\; \bigoplus_{\substack{Y \subset X \\ \text{prime divisor}}}
    \mathbb Z \cdot Y
  \]
  is the free abelian group on the set of prime divisors of \(X\)
  (\ref{def:codim1_cycles}); an element \(D = \sum n_i Y_i\) with finitely
  many nonzero coefficients \(n_i \in \mathbb Z\) is called a \emph{Weil
  divisor}, and is \emph{effective} when all \(n_i \geq 0\).
\end{definition}
\begin{definition}
\leanok
[Regular in codimension one]
  \label{def:isRegularInCodimensionOne}
  \lean{AlgebraicGeometry.Scheme.IsRegularInCodimensionOne}
  % SOURCE: [Hartshorne], II \S6, p.~130, the fourth clause of \((\ast)\)
  % (read from references/hartshorne-algebraic-geometry.pdf, PDF page 147).
  % Stacks pointers: tag 0BX5 (regular schemes), tag 02ME
  % (DVR-of-stalk characterisation at codim-1 points).
  % SOURCE QUOTE: "\((\ast)\) \(X\) is a noetherian integral separated scheme
  % which is regular in codimension one."
  \textit{Source: Hartshorne, II \S6, p.~130 (the fourth clause of \((\ast)\)).}
  A scheme \(X\) is \emph{regular in codimension one} when for every point
  \(x \in X\) with \(\mathrm{codim}_X \overline{\{x\}} = 1\) (equivalently,
  \(\texttt{Order.coheight}\,x = 1\)) the local ring \(\mathcal O_{X,x}\) is
  a discrete valuation ring.
\end{definition}
\begin{definition}

\leanok
[DVR stalk from regularity in codimension one]
  \label{def:isRegularInCodimOne_dvrStalk}
  \lean{AlgebraicGeometry.Scheme.IsRegularInCodimensionOne.instIsDiscreteValuationRingStalk}
  \uses{def:isRegularInCodimensionOne, def:prime_divisor}
  For an integral scheme \(X\) satisfying
  \(\texttt{IsRegularInCodimensionOne}\) (\ref{def:isRegularInCodimensionOne})
  and any prime divisor \(Y : X.\texttt{PrimeDivisor}\)
  (\ref{def:prime_divisor}), the stalk
  \(\mathcal O_{X, Y.\texttt{point}}\) is a discrete valuation ring,
  a direct consequence of the definition of \(\texttt{IsRegularInCodimensionOne}\).
\end{definition}
\begin{proof}
  Follows directly from the definition of
  \(\texttt{IsRegularInCodimensionOne}\) (\ref{def:isRegularInCodimensionOne}).
\end{proof}
\begin{definition}

\leanok
[Krull dimension \(\leq 1\) of a prime-divisor stalk]
  \label{def:isRegularInCodimOne_krullDimStalk}
  \lean{AlgebraicGeometry.Scheme.IsRegularInCodimensionOne.instKrullDimLEStalk}
  \uses{def:isRegularInCodimOne_dvrStalk, def:prime_divisor}
  For an integral scheme \(X\) satisfying
  \(\texttt{IsRegularInCodimensionOne}\) and any prime divisor
  \(Y : X.\texttt{PrimeDivisor}\), the stalk
  \(\mathcal O_{X, Y.\texttt{point}}\) has Krull dimension \(\leq 1\), obtained
  from the DVR structure (\ref{def:isRegularInCodimOne_dvrStalk}) via the chain
  ``DVR \(\Rightarrow\) principal ideal ring \(\Rightarrow\) Krull dimension
  \(\leq 1\)''.
\end{definition}
\begin{proof}
  Proved directly in Lean.
\end{proof}
\section{Open-immersion descent for prime divisors}
\label{sec:open_immersion_descent}
\textit{Source: Stacks Project tags 02IZ (stalk under open immersion) +
01J6 (open immersion preserves regularity); uses the project-local substrate
\texttt{AlgebraicJacobian/Albanese/CoheightBridge.lean}.}

The chapter's order / principal-divisor API is naturally stated on the
ambient scheme \(X\), but several downstream proofs (the
\(\texttt{rationalMap\_order\_finite\_support}\) non-zero branch in
particular) require localising to an affine open \(U \hookrightarrow X\),
identifying prime divisors of \(X\) lying in \(U\) with prime divisors of
\(U\), and transporting stalk-level data along the open-immersion stalk
isomorphism. This section records the substrate bijection
\(\{Y : X.\texttt{PrimeDivisor} \mid Y.\texttt{point} \in U\}
 \;\simeq\; U.\texttt{PrimeDivisor}\)
and the open-immersion descent of the \(\texttt{IsRegularInCodimensionOne}\)
class.

\begin{lemma}

\leanok
[Extensionality for prime divisors]
  \label{lem:primeDivisor_ext}
  \lean{AlgebraicGeometry.Scheme.PrimeDivisor.ext}
  \uses{def:prime_divisor}
  Two prime divisors \(Y, Y' : X.\texttt{PrimeDivisor}\)
  (\ref{def:prime_divisor}) on a scheme \(X\) are equal as soon as their
  underlying generic points agree: if \(Y.\texttt{point} = Y'.\texttt{point}\)
  then \(Y = Y'\). The codimension-one witness
  \(\texttt{Order.coheight}\,Y.\texttt{point} = 1\) is a \(\texttt{Prop}\)-valued
  field, hence proof-irrelevant, so it carries no data distinguishing two prime
  divisors that share the same underlying point. This is the
  \(\texttt{@[ext]}\)-generated extensionality lemma underlying the round-trip
  identities of the bijection \(\texttt{equivOpen}\)
  (\ref{lem:primeDivisor_equivOpen}).
\end{lemma}
\begin{proof}
  Proved directly in Lean.
\end{proof}

\begin{definition}

\leanok
[Restriction of a prime divisor to an open subscheme]
  \label{def:primeDivisor_restrictToOpen}
  \lean{AlgebraicGeometry.Scheme.PrimeDivisor.restrictToOpen}
  \uses{def:prime_divisor}
  Given a scheme \(X\), an open subscheme \(U \hookrightarrow X\), a prime
  divisor \(Y : X.\texttt{PrimeDivisor}\), and a witness \(Y.\texttt{point} \in U\),
  the restriction \(Y |_U : U.\texttt{toScheme}.\texttt{PrimeDivisor}\) packages
  the point \(Y.\texttt{point}\) (viewed as a point of \(U\) via the open
  inclusion) with the coheight witness from the project-local substrate
  \(\texttt{Order.coheight\_eq\_of\_isOpenEmbedding}\).
\end{definition}
\begin{definition}

\leanok
[Push of a prime divisor along an open immersion]
  \label{def:primeDivisor_ofOpen}
  \lean{AlgebraicGeometry.Scheme.PrimeDivisor.ofOpen}
  \uses{def:prime_divisor}
  Given a scheme \(X\), an open subscheme \(U \hookrightarrow X\), and a
  prime divisor \(Y : U.\texttt{toScheme}.\texttt{PrimeDivisor}\), the push
  \(\texttt{ofOpen}\,U\,Y : X.\texttt{PrimeDivisor}\) packages the point
  \(Y.\texttt{point} \in U \hookrightarrow X\) (viewed as a point of \(X\))
  with the coheight witness from
  \(\texttt{Order.coheight\_eq\_of\_isOpenEmbedding}\) applied in the
  opposite direction.
\end{definition}
\begin{lemma}

\leanok
[Bijection of prime divisors under open immersion]
  \label{lem:primeDivisor_equivOpen}
  \lean{AlgebraicGeometry.Scheme.PrimeDivisor.equivOpen}
  \uses{def:primeDivisor_restrictToOpen, def:primeDivisor_ofOpen}
  For an open subscheme \(U \hookrightarrow X\) of a scheme \(X\), there is a
  bijection
  \[
    \{Y : X.\texttt{PrimeDivisor} \mid Y.\texttt{point} \in U\}
    \;\simeq\;
    U.\texttt{toScheme}.\texttt{PrimeDivisor}
  \]
  packaging the restriction (\ref{def:primeDivisor_restrictToOpen}) and the
  push (\ref{def:primeDivisor_ofOpen}) as mutual inverses. Both
  \texttt{left\_inv} and \texttt{right\_inv} hold definitionally
  (\texttt{rfl}) because both directions preserve the underlying point
  field via subtype eta + structure eta.
\end{lemma}
\begin{lemma}

\leanok
[Point field of a restricted prime divisor]
  \label{lem:primeDivisor_restrictToOpen_point}
  \lean{AlgebraicGeometry.Scheme.PrimeDivisor.restrictToOpen_point}
  \uses{def:primeDivisor_restrictToOpen}
  For an open subscheme \(U \hookrightarrow X\), a prime divisor
  \(Y : X.\texttt{PrimeDivisor}\), and a witness \(h : Y.\texttt{point} \in U\),
  the underlying point of the restriction
  \(\texttt{restrictToOpen}\,U\,Y\,h\) (\ref{def:primeDivisor_restrictToOpen})
  is the pair \(\langle Y.\texttt{point}, h\rangle \in U\); i.e.\ restriction
  preserves the underlying point field.
\end{lemma}
\begin{proof}
  Proved directly in Lean (the equality is definitional).
\end{proof}
\begin{lemma}

\leanok
[Point field of a pushed prime divisor]
  \label{lem:primeDivisor_ofOpen_point}
  \lean{AlgebraicGeometry.Scheme.PrimeDivisor.ofOpen_point}
  \uses{def:primeDivisor_ofOpen}
  For an open subscheme \(U \hookrightarrow X\) and a prime divisor
  \(Y_U : U.\texttt{toScheme}.\texttt{PrimeDivisor}\), the underlying point of
  the push \(\texttt{ofOpen}\,U\,Y_U\) (\ref{def:primeDivisor_ofOpen}) is the
  image \(Y_U.\texttt{point} \in U \hookrightarrow X\); i.e.\ the push
  preserves the underlying point field.
\end{lemma}
\begin{proof}
  Proved directly in Lean (the equality is definitional).
\end{proof}
\begin{lemma}

\leanok
[Stalk isomorphism across an open immersion]
  \label{lem:primeDivisor_stalkIso}
  \lean{AlgebraicGeometry.Scheme.PrimeDivisor.stalkIso}
  \uses{def:primeDivisor_restrictToOpen}
  \textit{Mathlib-backed: a thin wrapper around
  \texttt{AlgebraicGeometry.Scheme.Opens.stalkIso} (stalk invariance under
  restriction to a non-empty open neighbourhood).}
  For \(Y : X.\texttt{PrimeDivisor}\) with \(Y.\texttt{point} \in U\), the
  stalk presheaf of \(U.\texttt{toScheme}\) at the restricted point coincides,
  up to canonical isomorphism, with the stalk presheaf of \(X\) at
  \(Y.\texttt{point}\). The Lean wrapper packages Mathlib's
  \(\texttt{Scheme.Opens.stalkIso}\) for the project's
  \texttt{PrimeDivisor} bundle.
\end{lemma}
\begin{theorem}

\leanok
[\(\texttt{IsRegularInCodimensionOne}\) descends along open immersion]
  \label{thm:isRegularInCodimensionOne_open}
  \lean{AlgebraicGeometry.Scheme.IsRegularInCodimensionOne.instOpen}
  \uses{def:isRegularInCodimensionOne, lem:primeDivisor_stalkIso,
    def:primeDivisor_ofOpen}
  Let \(X\) be a scheme satisfying \(\texttt{IsRegularInCodimensionOne}\)
  and \(U \hookrightarrow X\) an open subscheme with
  \(\texttt{IsIntegral}\,U.\texttt{toScheme}\) (the latter is \emph{not}
  automatic from integrality of \(X\); it is threaded explicitly).
  Then \(U.\texttt{toScheme}\) also satisfies
  \(\texttt{IsRegularInCodimensionOne}\); the DVR-of-stalk property at a
  codim-1 point \(Y : U.\texttt{toScheme}.\texttt{PrimeDivisor}\) is
  transported from the corresponding point
  \(\texttt{ofOpen}\,U\,Y : X.\texttt{PrimeDivisor}\) on the ambient
  scheme via the stalk isomorphism (\ref{lem:primeDivisor_stalkIso})
  together with
  \(\texttt{IsDiscreteValuationRing.RingEquivClass.isDiscreteValuationRing}\)
  on the ring-equiv image.
\end{theorem}

\subsection{Naturality of the order-of-vanishing across a ring isomorphism}
\label{subsec:ordFrac_ringEquiv}
\textit{Source: Project-bespoke Mathlib supplement, built on Mathlib's
\texttt{Ring.ord} / \texttt{Ring.ordMonoidWithZeroHom} / \texttt{Ring.ordFrac}
order-of-vanishing API (\texttt{Mathlib.RingTheory.OrderOfVanishing.Basic},
Stacks tag 02MD).}

Mathlib's order-of-vanishing data --- the length valuation \(\texttt{ord}\), its
monoid-with-zero extension \(\texttt{ordMonoidWithZeroHom}\), and the
fraction-field hom \(\texttt{ordFrac}\) --- is natural across a ring
isomorphism.

\begin{lemma}

\leanok
[Order length is invariant under a ring isomorphism]
  \label{lem:ord_ringEquiv}
  \lean{Ring.ord_ringEquiv}
  For a ring isomorphism \(e \colon R \simeq S\) of commutative rings and an
  element \(x \in R\), the length valuations agree:
  \(\texttt{ord}\,R\,x = \texttt{ord}\,S\,(e\,x)\). Equivalently, the
  \(R\)-module length of \(R/(x)\) equals the \(S\)-module length of
  \(S/(e\,x)\), transported across \(e\).
\end{lemma}
\begin{proof}
  Proved directly in Lean: the quotient-ring isomorphism
  \(R/(x) \simeq S/(e\,x)\) induced by \(e\) is an \(R\)-linear equivalence
  (with the \(S\)-side carrying its \(R\)-module structure via \(e\)), and
  module length is preserved by surjective semilinear maps.
\end{proof}
\begin{lemma}

\leanok
[Non-zero divisors are preserved by a ring isomorphism]
  \label{lem:nonZeroDivisors_ringEquiv}
  \lean{Ring.nonZeroDivisors_ringEquiv}
  For a ring isomorphism \(e \colon R \simeq S\) of commutative rings and
  \(r \in R\), one has \(r \in R^{\times}_{\mathrm{nzd}}\) if and only if
  \(e\,r \in S^{\times}_{\mathrm{nzd}}\), where \((-)^{\times}_{\mathrm{nzd}}\)
  denotes the submonoid of non-zero divisors. A convenience \(\mathrm{Iff}\)
  repackaging of Mathlib's \(\texttt{MulEquivClass.map\_nonZeroDivisors}\).
\end{lemma}
\begin{proof}
  Proved directly in Lean from \(\texttt{MulEquivClass.map\_nonZeroDivisors}\)
  together with injectivity of \(e\).
\end{proof}
\begin{lemma}

\leanok
[The order monoid-with-zero hom is invariant under a ring isomorphism]
  \label{lem:ordMonoidWithZeroHom_ringEquiv}
  \lean{Ring.ordMonoidWithZeroHom_ringEquiv}
  \uses{lem:ord_ringEquiv, lem:nonZeroDivisors_ringEquiv}
  For a ring isomorphism \(e \colon R \simeq S\) of nontrivial commutative
  rings and \(r \in R\),
  \(\texttt{ordMonoidWithZeroHom}\,S\,(e\,r)
   = \texttt{ordMonoidWithZeroHom}\,R\,r\). The proof case-splits on whether
  \(r\) is a non-zero divisor (\ref{lem:nonZeroDivisors_ringEquiv}) and uses
  \(\texttt{ord}\) invariance (\ref{lem:ord_ringEquiv}) on the non-zero-divisor
  branch.
\end{lemma}
\begin{proof}
  Proved directly in Lean.
\end{proof}
\begin{lemma}

\leanok
[The fraction-field order-of-vanishing is invariant under a ring isomorphism]
  \label{lem:ordFrac_ringEquiv}
  \lean{Ring.ordFrac_ringEquiv}
  \uses{lem:ordMonoidWithZeroHom_ringEquiv, lem:nonZeroDivisors_ringEquiv}
  Let \(e \colon R \simeq S\) be a ring isomorphism between Noetherian
  integral domains of Krull dimension \(\leq 1\), and let
  \(e_K \colon K_R \simeq K_S\) be a compatible isomorphism of their fraction
  fields (\(e_K(\iota_R\,r) = \iota_S(e\,r)\) for \(r \in R\), where
  \(\iota\) denotes the canonical inclusion into the fraction field). Then the
  order-of-vanishing hom on fraction fields is invariant:
  \(\texttt{ordFrac}\,S\,(e_K\,x) = \texttt{ordFrac}\,R\,x\) for every
  \(x \in K_R\).
\end{lemma}
\begin{proof}
  Proved directly in Lean. The junk-zero case \(x = 0\) is immediate; for
  \(x \neq 0\), write \(x = a/b\) with \(a, b\) non-zero divisors via
  \(\texttt{IsLocalization.surj}\), so \(e_K\,x = (e\,a)/(e\,b)\) by the
  compatibility hypothesis; both sides reduce via
  \(\texttt{Ring.ordFrac\_eq\_div}\) to a quotient of
  \(\texttt{ordMonoidWithZeroHom}\)-values, which agree numerator- and
  denominator-wise by \ref{lem:ordMonoidWithZeroHom_ringEquiv}.
\end{proof}
\begin{definition}\leanok
[Function-field isomorphism along an open immersion]
  \label{def:functionFieldIso}
  \lean{AlgebraicGeometry.Scheme.Opens.functionFieldIso}
  \uses{lem:primeDivisor_stalkIso}
  \textit{Source: Archon-original packaging; built on Mathlib's
  \texttt{AlgebraicGeometry.Scheme.Opens.stalkIso} +
  \texttt{genericPoint\_eq\_of\_isOpenImmersion}.}
  For an integral scheme \(X\) and a non-empty open subscheme
  \(U \hookrightarrow X\) with \(\texttt{IsIntegral}\,U.\texttt{toScheme}\),
  the function field of \(U.\texttt{toScheme}\) is canonically isomorphic
  to the function field of \(X\). The isomorphism is built by composing
  \(\texttt{Scheme.Opens.stalkIso}\,U\,(\texttt{genericPoint}\,U.\texttt{toScheme})\)
  with \(X.\texttt{presheaf.stalkCongr}\) along the equality
  \((\texttt{genericPoint}\,U.\texttt{toScheme}).\texttt{val}
   = \texttt{genericPoint}\,X\)
  coming from \(\texttt{genericPoint\_eq\_of\_isOpenImmersion}\,U.\iota\).
  This is the entry point that discharges the \(\texttt{h\_compat}\)
  hypothesis of
  \(\texttt{Scheme.PrimeDivisor.ordFrac\_stalkIso\_naturality}\),
  promoting the parameterised naturality to the canonical function-field
  iso of integral schemes.
\end{definition}


\begin{lemma}\leanok
[Function-field-iso compatibility, morphism level]
  \label{lem:functionFieldIso_compat}
  \lean{AlgebraicGeometry.Scheme.PrimeDivisor.functionFieldIso_compat}
  \uses{def:functionFieldIso, lem:primeDivisor_stalkIso}
  For an integral scheme \(X\), a nonempty integral open \(U\), and a
  prime divisor \(Y\) with \(Y.\texttt{point} \in U\), the square
  \(\texttt{stalkIso}\,U\,Y \,\ggg\, \texttt{stalkSpecializes}\,
   (\text{gp}_X \rightsquigarrow Y.\texttt{point})
   = \texttt{stalkSpecializes}\,(\text{gp}_U \rightsquigarrow Y')
     \,\ggg\, \texttt{functionFieldIso}\,U\)
  commutes in \(\texttt{CommRingCat}\), where the horizontal maps are
  the canonical algebra maps \(\mathcal O_{\cdot, Y} \to K(\cdot)\) to
  the respective generic points.
\end{lemma}
\begin{proof}
  \uses{def:functionFieldIso}
  A germ-universal chase via \(\texttt{stalk\_hom\_ext}\): after fixing
  an open \(V\) containing the point, both composites reduce to
  \(\texttt{germ} \,\ggg\, \texttt{stalkSpecializes}\) and agree by
  \(\texttt{germ\_stalkSpecializes}\). The identification of
  \(\texttt{stalkCongr}\) with \(\texttt{stalkSpecializes}\) along
  inseparability is definitional.
\end{proof}

\begin{lemma}

\leanok
[Parameterised \(\texttt{ordFrac}\) naturality across the open-immersion stalk iso]
  \label{lem:ordFrac_stalkIso_naturality}
  \lean{AlgebraicGeometry.Scheme.PrimeDivisor.ordFrac_stalkIso_naturality}
  \uses{lem:primeDivisor_stalkIso, lem:ordFrac_ringEquiv,
    def:primeDivisor_restrictToOpen, def:isRegularInCodimensionOne}
  Let \(X\) be an integral, locally Noetherian scheme satisfying
  \(\texttt{IsRegularInCodimensionOne}\), let \(U \hookrightarrow X\) be a
  nonempty integral open subscheme (also regular in codimension one), and let
  \(Y : X.\texttt{PrimeDivisor}\) with \(Y.\texttt{point} \in U\). Given any
  fraction-field isomorphism
  \(e_K \colon U.\texttt{toScheme}.\texttt{functionField} \simeq
   X.\texttt{functionField}\) compatible with the stalk isomorphism
  \(\texttt{stalkIso}\,U\,Y\) (\ref{lem:primeDivisor_stalkIso}) over the local
  rings, the order-of-vanishing on the two function fields agrees:
  \(\texttt{ordFrac}\,(\mathcal O_{X, Y.\texttt{point}})\,(e_K\,f)
   = \texttt{ordFrac}\,(\mathcal O_{U, Y'})\,f\)
  for every \(f \in U.\texttt{toScheme}.\texttt{functionField}\), where
  \(Y' = (\texttt{restrictToOpen}\,U\,Y).\texttt{point}\). This is the
  scheme-level packaging of \ref{lem:ordFrac_ringEquiv}, with the ring iso
  taken to be the stalk iso and the compatibility hypothesis left as a
  parameter (discharged in \ref{lem:order_eq_order_restrict}).
\end{lemma}
\begin{proof}
  \uses{lem:ordFrac_ringEquiv}
  Proved directly in Lean: apply \ref{lem:ordFrac_ringEquiv} to the ring iso
  \(\texttt{stalkIso}\,U\,Y\) (a ring isomorphism of Noetherian
  Krull-dimension-\(\leq 1\) domains, the stalks being DVRs by
  \ref{def:isRegularInCodimensionOne}) and the supplied compatible
  fraction-field iso \(e_K\).
\end{proof}
\begin{lemma}\leanok
[Order is invariant under restriction to an open chart]
  \label{lem:order_eq_order_restrict}
  \lean{AlgebraicGeometry.Scheme.PrimeDivisor.order_eq_order_restrict}
  \uses{lem:functionFieldIso_compat, lem:ordFrac_stalkIso_naturality}
  Under \([\texttt{IsIntegral}\,X]\,[\texttt{IsLocallyNoetherian}\,X]\,
  [\texttt{IsRegularInCodimensionOne}\,X]\), for a nonempty integral
  open \(U\), a prime divisor \(Y\) with \(Y.\texttt{point} \in U\),
  and \(f \in U.\texttt{toScheme}.\texttt{functionField}\),
  \(\texttt{order}\,Y\,(\texttt{functionFieldIso}\,U\,f)
   = \texttt{order}\,(\texttt{restrictToOpen}\,U\,Y)\,f\).
\end{lemma}
\begin{proof}
  \uses{lem:functionFieldIso_compat}
  Discharge the \(\texttt{h\_compat}\) hypothesis of
  \(\texttt{ordFrac\_stalkIso\_naturality}\) from
  \cref{lem:functionFieldIso_compat} (the pointwise form follows by
  \(\texttt{congrArg}\) and is definitionally equal to the required
  algebra-map factorisation), then unfold \(\texttt{order}\) and
  rewrite by the naturality lemma.
\end{proof}

\section{Order of a rational function at a prime divisor}
\begin{definition}
\leanok
[Order of a rational function at a prime divisor]
  \label{def:order_at_point}

  \uses{def:codim1_cycles, def:isRegularInCodimOne_krullDimStalk}
  \lean{AlgebraicGeometry.Scheme.RationalMap.order}
  % SOURCE: [Hartshorne], II.6, p.~130--131 (valuation \(v_Y\) of a prime
  % divisor; "zero/pole of \(f\) along \(Y\)") (read from
  % references/hartshorne-algebraic-geometry.pdf, PDF pages 147--148).
  % Stacks pointer: tag 02ME.
  % SOURCE QUOTE: "If \(Y\) is a prime divisor on \(X\), let \(\eta \in Y\) be its
  % generic point. Then the local ring \(\mathcal O_{\eta, X}\) of \(X\) at
  % \(\eta\) is a discrete valuation ring with quotient field \(K\), the
  % function field of \(X\). We call the corresponding discrete valuation
  % \(v_Y\) the \emph{valuation of \(Y\)}."
  % SOURCE QUOTE: "Now let \(f \in K^\ast\) be any nonzero rational function
  % on \(X\). Then \(v_Y(f)\) is an integer. If it is positive, we say \(f\) has
  % a \emph{zero} along \(Y\), of that order; if it is negative, we say \(f\)
  % has a \emph{pole} along \(Y\), of order \(-v_Y(f)\)."
  \textit{Source: Hartshorne, II.6, pp.~130--131 (valuation \(v_Y\)).}
  Let \(X\) satisfy \((*)\). For a prime divisor \(Y \subset X\) with generic
  point \(\eta \in Y\), the local ring \(\mathcal O_{X, \eta}\) is a discrete
  valuation ring (by regularity in codimension one) with quotient field the
  function field \(K(X)\). For a nonzero rational function \(f \in K(X)^{\times}\)
  the \emph{order of \(f\) along \(Y\)} is the integer
  \[
    \ord_Y(f) \;:=\; v_Y(f) \;\in\; \mathbb Z,
  \]
  where \(v_Y\) is the normalised discrete valuation associated to
  \(\mathcal O_{X, \eta}\). If \(\ord_Y(f) > 0\), \(f\) has a \emph{zero of
  order \(\ord_Y(f)\) along \(Y\)}; if \(\ord_Y(f) < 0\), \(f\) has a \emph{pole
  of order \(-\ord_Y(f)\) along \(Y\)}.
  Specialisation to a curve. If \(C\) is a smooth proper curve over \(\bar k\)
  and \(P \in C\) is a closed point, then \(P\) is automatically a prime
  divisor (it is closed, integral, of codimension one in the \(1\)-dimensional
  scheme \(C\)); the local ring \(\mathcal O_{C, P}\) is a discrete valuation
  ring by smoothness, and \(\ord_P(f) = v_P(f)\) is the standard DVR
  valuation of \(f\) at \(P\).
\end{definition}
\subsection{Valuation identities for the order function}
\label{subsec:order_identities}
The order function \(\ord_Y\) inherits the homomorphism laws of the
underlying discrete valuation (with Mathlib's junk-on-zero convention). The
following identities feed the principal-divisor homomorphism
(\ref{thm:principal_hom}) and the ramification--inertia computation in
\ref{lem:degree_positivePart_principal_eq_finrank}.
\begin{lemma}
[Order of zero]
  \label{lem:order_zero}
  \lean{AlgebraicGeometry.Scheme.RationalMap.order_zero}
  \uses{def:order_at_point}
  For a prime divisor \(Y\) of \(X\), \(\ord_Y(0) = 0\), the junk-on-zero
  value of \ref{def:order_at_point}.
\end{lemma}
\begin{proof}
  Proved directly in Lean (\(\texttt{ordFrac}\) is a monoid-with-zero hom,
  so it sends \(0\) to \(0\), and \(\texttt{WithZero.log}\,0 = 0\)).
\end{proof}
\begin{lemma}
[Order of one]
  \label{lem:order_one}
  \lean{AlgebraicGeometry.Scheme.RationalMap.order_one}
  \uses{def:order_at_point}
  For a prime divisor \(Y\) of \(X\), \(\ord_Y(1) = 0\).
\end{lemma}
\begin{proof}
  Proved directly in Lean (from \(\texttt{map\_one}\) of \(\texttt{ordFrac}\)
  and \(\texttt{WithZero.log\_one}\)).
\end{proof}
\begin{lemma}
[Order of a product of nonzero functions]
  \label{lem:order_mul_of_ne_zero}
  \lean{AlgebraicGeometry.Scheme.RationalMap.order_mul_of_ne_zero}
  \uses{def:order_at_point}
  For a prime divisor \(Y\) of \(X\) and nonzero
  \(f, g \in K(X)^{\times}\),
  \(\ord_Y(fg) = \ord_Y(f) + \ord_Y(g)\).
\end{lemma}
\begin{proof}
  Proved directly in Lean (from \(\texttt{map\_mul}\) of \(\texttt{ordFrac}\)
  and \(\texttt{WithZero.log\_mul}\), which needs both factors nonzero).
\end{proof}
\begin{lemma}
[Order of an inverse]
  \label{lem:order_inv}
  \lean{AlgebraicGeometry.Scheme.RationalMap.order_inv}
  \uses{def:order_at_point}
  For a prime divisor \(Y\) of \(X\) and any \(f \in K(X)\),
  \(\ord_Y(f^{-1}) = -\ord_Y(f)\) (the \(f = 0\) case is junk-vacuous, both
  sides being zero).
\end{lemma}
\begin{proof}
  Proved directly in Lean (from \(\texttt{map\_inv}_0\) of \(\texttt{ordFrac}\)
  and \(\texttt{WithZero.log\_inv}\)).
\end{proof}
\begin{lemma}
[Order of a unit inverse]
  \label{lem:order_units_inv}
  \lean{AlgebraicGeometry.Scheme.RationalMap.order_units_inv}
  \uses{lem:order_inv}
  For a prime divisor \(Y\) of \(X\) and a unit \(u \in K(X)^{\times}\),
  \(\ord_Y(u^{-1}) = -\ord_Y(u)\). A unit-typed specialisation of
  \ref{lem:order_inv}, convenient at the call sites that thread the
  multiplicative group \(K(X)^{\times}\) (e.g.\ \ref{thm:principal_hom}).
\end{lemma}
\begin{proof}
  Proved directly in Lean (rewrite the unit inverse to the field inverse and
  apply \ref{lem:order_inv}).
\end{proof}
\begin{lemma}
[Order of a negation]
  \label{lem:order_neg}
  \lean{AlgebraicGeometry.Scheme.RationalMap.order_neg}
  \uses{def:order_at_point}
  For a prime divisor \(Y\) of \(X\) and any \(f \in K(X)\),
  \(\ord_Y(-f) = \ord_Y(f)\).
\end{lemma}
\begin{proof}
  Proved directly in Lean: \((-f)^2 = f^2\), so applying \(\texttt{ordFrac}\)
  and \(\texttt{WithZero.log\_pow}\) gives \(2\,\ord_Y(-f) = 2\,\ord_Y(f)\) in
  \(\mathbb Z\); cancelling \(2 \neq 0\) yields the equality.
\end{proof}
\begin{lemma}
[Order of a power of a nonzero function]
  \label{lem:order_pow_of_ne_zero}
  \lean{AlgebraicGeometry.Scheme.RationalMap.order_pow_of_ne_zero}
  \uses{lem:order_mul_of_ne_zero, lem:order_one}
  For a prime divisor \(Y\) of \(X\), a nonzero \(f \in K(X)^{\times}\), and a
  natural number \(n\), \(\ord_Y(f^n) = n \cdot \ord_Y(f)\).
\end{lemma}
\begin{proof}
  Proved directly in Lean by induction on \(n\) from
  \ref{lem:order_mul_of_ne_zero} and \ref{lem:order_one}.
\end{proof}
\section{Divisor of a closed point on a curve}
\begin{definition}
\leanok
[Divisor of a closed point on a curve]
  \label{def:divisor_closed_point}
  \uses{def:codim1_cycles}
  
  \lean{AlgebraicGeometry.Scheme.WeilDivisor.ofClosedPoint}
  % SOURCE: [Hartshorne], II.6 ("Divisors on Curves"), p.~137 (read from
  % references/hartshorne-algebraic-geometry.pdf, PDF page 154):
  % "A prime divisor is just a closed point, so an arbitrary divisor can
  % be written \(D = \sum n_i P_i\), where the \(P_i\) are closed points, and
  % \(n_i \in \mathbf Z\)."
  % SOURCE QUOTE: "Now we come to the study of divisors on curves. If \(X\)
  % is a nonsingular curve, then \(X\) satisfies the condition \((\ast)\) used
  % above, so we can talk about divisors on \(X\). A prime divisor is just
  % a closed point, so an arbitrary divisor can be written $D = \sum n_i
  % P_i\(, where the \)P_i\( are closed points, and \)n_i \in \mathbf Z$."
  \textit{Source: Hartshorne, II.6, p.~137 (``Divisors on Curves'').}
  Let \(C\) be a smooth proper curve over a field \(k\) and let \(P \in C\) be a
  closed point. The \emph{Weil divisor associated to \(P\)} is the element
  \[
    [P] \;:=\; 1 \cdot P \;\in\; \mathrm{Div}(C)
  \]
  i.e.\ the prime divisor \(P\) with multiplicity one. The assignment is
  well-defined because the closed point \(P\) is automatically a codimension-\(1\)
  integral closed subscheme of the \(1\)-dimensional integral scheme \(C\).
  Conversely, every prime divisor on \(C\) is of this form, so an arbitrary
  divisor on \(C\) is a finite formal \(\mathbb Z\)-linear combination
  \(\sum n_i [P_i]\) of closed points.
\end{definition}
\begin{lemma}
\leanok
[Bridge equation: \texttt{ofClosedPoint} in the codim-\(1\) regime]
  \label{lem:ofClosedPoint_eq_single}
  \lean{AlgebraicGeometry.Scheme.WeilDivisor.ofClosedPoint_eq_single}
  \uses{def:divisor_closed_point}
  Let \(C : \texttt{Scheme.\{u\}}\) be a scheme and let \(P \in C\) be a
  point with \(\texttt{IsClosed}\,(\{P\} : \texttt{Set}\ C)\). If
  \(h : \texttt{Order.coheight}\ P = 1\) (the codimension-one witness
  promoting \(P\) to a \texttt{PrimeDivisor} via \ref{def:divisor_closed_point}),
  then
  \[
    \texttt{ofClosedPoint}\ P\ \_\!hP \;=\;
      \texttt{Finsupp.single}\ \langle P, h\rangle\ 1
      \;\in\; \mathrm{Div}(C).
  \]
\end{lemma}
\begin{proof}
  Unfolds the dependent-\texttt{if} in the definition of
  \texttt{ofClosedPoint} (\ref{def:divisor_closed_point}) on the
  positive branch using the hypothesis \(h\).
\end{proof}
\begin{lemma}
\leanok
[Bridge equation: \texttt{ofClosedPoint} junk-branch]
  \label{lem:ofClosedPoint_eq_zero}
  \lean{AlgebraicGeometry.Scheme.WeilDivisor.ofClosedPoint_eq_zero}
  \uses{def:divisor_closed_point}
  Let \(C : \texttt{Scheme.\{u\}}\) be a scheme and let \(P \in C\) be a
  point with \(\texttt{IsClosed}\,(\{P\} : \texttt{Set}\ C)\). If
  \(\texttt{Order.coheight}\ P \neq 1\) (the off-branch: \(P\) is not a
  codimension-one point of \(C\) --- e.g.\ \(P\) is the generic point of an
  integral scheme, or a higher-codim point of a higher-dimensional
  scheme), then
  \[
    \texttt{ofClosedPoint}\ P\ \_\!hP \;=\;
      0 \;\in\; \mathrm{Div}(C).
  \]
\end{lemma}
\begin{proof}
  Unfolds the dependent-\texttt{if} in the definition of
  \texttt{ofClosedPoint} (\ref{def:divisor_closed_point}) on the
  negative branch using the hypothesis
  \(\texttt{Order.coheight}\ P \neq 1\).
\end{proof}
\section{Degree of a divisor over an algebraically closed base}
\begin{definition}
\leanok
[Degree of a divisor on a curve over \(\bar k\)]
  \uses{def:codim1_cycles}
  \label{def:divisor_degree}
  
  \lean{AlgebraicGeometry.Scheme.WeilDivisor.degree}
  % SOURCE: [Hartshorne], II.6 ("Divisors on Curves"), p.~137 (read from
  % references/hartshorne-algebraic-geometry.pdf, PDF page 154).
  % Stacks pointer: tag 0BE0.
  % SOURCE QUOTE: "We define the \emph{degree} of \(D\) to be \(\sum n_i\)."
  \textit{Source: Hartshorne, II.6, p.~137 (definition of degree on a
  nonsingular curve).}
  Let \(C\) be a smooth proper curve over an algebraically closed field
  \(\bar k\). The \emph{degree} of a Weil divisor $D = \sum n_i [P_i] \in
  \mathrm{Div}(C)$ is the integer
  \[
    \deg(D) \;:=\; \sum_i n_i \;\in\; \mathbb Z.
  \]
  The definition uses that every closed point of \(C\) has residue field
  \(\bar k\) (so each prime divisor contributes degree one); over a general
  field \(k\) the same definition with the residue-field-degree weight
  \(\deg(D) := \sum_i n_i \, [\kappa(P_i) : k]\) recovers the
  \emph{geometric} degree, but the project's RR bridge needs only the
  \(\bar k\) specialisation.
\end{definition}
\begin{theorem}
\leanok
[The degree map is a group homomorphism]
  \label{thm:divisor_degree_hom}
  \lean{AlgebraicGeometry.Scheme.WeilDivisor.degree_hom}
  \uses{def:divisor_degree, def:codim1_cycles}
  Let \(C\) be a smooth proper curve over an algebraically closed field
  \(\bar k\). The degree map of \ref{def:divisor_degree} is a group
  homomorphism
  \[
    \deg \colon \mathrm{Div}(C) \longrightarrow \mathbb Z,
    \qquad
    \deg(D_1 + D_2) \;=\; \deg(D_1) + \deg(D_2),
    \qquad
    \deg(-D) \;=\; -\deg(D).
  \]
\end{theorem}
\begin{proof}
  Immediate from the definition: \(\mathrm{Div}(C)\) is the free abelian group on the
  set of closed points of \(C\), and \(\deg\) is the unique group homomorphism
  sending each generator \([P]\) to \(1 \in \mathbb Z\) (equivalently: the sum
  of the multiplicities). Both additivity and negation thus reduce to the
  corresponding identities for finite sums of integers.
\end{proof}
\begin{lemma}

\leanok
[Degree homomorphism agrees with the degree map]
  \label{lem:degree_hom_apply}
  \lean{AlgebraicGeometry.Scheme.WeilDivisor.degree_hom_apply}
  \uses{thm:divisor_degree_hom, def:divisor_degree}
  For every Weil divisor \(D \in \mathrm{Div}(X)\), the bundled group
  homomorphism \(\deg_{\mathrm{hom}}\) (\ref{thm:divisor_degree_hom}) applied
  to \(D\) equals the degree \(\deg(D)\) (\ref{def:divisor_degree}). This is
  the unfolding identity tying the \texttt{AddMonoidHom} packaging to the bare
  sum-of-coefficients.
\end{lemma}
\begin{proof}
  Proved directly in Lean (from \(\texttt{Finsupp.liftAddHom\_apply}\)).
\end{proof}
\begin{lemma}

\leanok
[Degree of the zero divisor]
  \label{lem:degree_zero}
  \lean{AlgebraicGeometry.Scheme.WeilDivisor.degree_zero}
  \uses{lem:degree_hom_apply, def:divisor_degree}
  \(\deg(0) = 0 \in \mathbb Z\).
\end{lemma}
\begin{proof}
  Proved directly in Lean (\(\deg_{\mathrm{hom}}\) is an \texttt{AddMonoidHom}
  so sends \(0\) to \(0\); use \ref{lem:degree_hom_apply}).
\end{proof}
\begin{lemma}

\leanok
[Degree is additive]
  \label{lem:degree_add}
  \lean{AlgebraicGeometry.Scheme.WeilDivisor.degree_add}
  \uses{lem:degree_hom_apply, def:divisor_degree}
  For Weil divisors \(D_1, D_2 \in \mathrm{Div}(X)\),
  \(\deg(D_1 + D_2) = \deg(D_1) + \deg(D_2)\).
\end{lemma}
\begin{proof}
  Proved directly in Lean (additivity of \(\deg_{\mathrm{hom}}\) via
  \ref{lem:degree_hom_apply}).
\end{proof}
\begin{lemma}

\leanok
[Degree of a negation]
  \label{lem:degree_neg}
  \lean{AlgebraicGeometry.Scheme.WeilDivisor.degree_neg}
  \uses{lem:degree_hom_apply, def:divisor_degree}
  For a Weil divisor \(D \in \mathrm{Div}(X)\), \(\deg(-D) = -\deg(D)\).
\end{lemma}
\begin{proof}
  Proved directly in Lean (\(\deg_{\mathrm{hom}}\) preserves negation; use
  \ref{lem:degree_hom_apply}).
\end{proof}
\begin{lemma}

\leanok
[Degree is subtractive]
  \label{lem:degree_sub}
  \lean{AlgebraicGeometry.Scheme.WeilDivisor.degree_sub}
  \uses{lem:degree_hom_apply, def:divisor_degree}
  For Weil divisors \(D_1, D_2 \in \mathrm{Div}(X)\),
  \(\deg(D_1 - D_2) = \deg(D_1) - \deg(D_2)\).
\end{lemma}
\begin{proof}
  Proved directly in Lean (subtractivity of \(\deg_{\mathrm{hom}}\) via
  \ref{lem:degree_hom_apply}).
\end{proof}
\begin{lemma}

\leanok
[Degree of a one-point divisor]
  \label{lem:degree_single}
  \lean{AlgebraicGeometry.Scheme.WeilDivisor.degree_single}
  \uses{def:divisor_degree, def:prime_divisor}
  For a prime divisor \(Y\) and \(n \in \mathbb Z\), the degree of the
  one-point-supported Weil divisor \(n \cdot Y\) is \(n\):
  \(\deg(n \cdot Y) = n\).
\end{lemma}
\begin{proof}
  Proved directly in Lean (from \(\texttt{Finsupp.sum\_single\_index}\)).
\end{proof}
\section{Principal divisors}
\begin{lemma}

\leanok
[Finite support of the order function — Hartshorne 6.1 / Stacks 02RV]
  \label{lem:rationalMap_order_finite_support}
  \lean{AlgebraicGeometry.rationalMap_order_finite_support}
  \uses{def:order_at_point, def:prime_divisor}
  % SOURCE: [Hartshorne], II.6, Lemma~6.1, p.~131 (read from
  % references/hartshorne-algebraic-geometry.pdf, PDF page 148).
  % SOURCE QUOTE: "Lemma 6.1. Let \(X\) satisfy \((\ast)\), and let \(f \in K^\ast\)
  % be a nonzero function on \(X\). Then \(v_Y(f) = 0\) for all except finitely
  % many prime divisors \(Y\)."
  \textit{Source: Hartshorne, II.6, Lemma~6.1, p.~131 (a.k.a.\ Stacks 02RV).}

  Let \(X\) be an integral scheme that is regular in codimension one and is
  \emph{Noetherian} --- equivalently, both \emph{locally Noetherian} and
  \emph{quasi-compact}, i.e.\ admitting a finite affine cover. Let
  \(f \in K(X)\) be a nonzero rational function. Then
  \(\{Y \in X^{(1)} \mid \ord_Y(f) \ne 0\}\) is finite. Both hypotheses are
  required: local Noetherianness makes each affine chart-ring Noetherian, and
  quasi-compactness bounds the number of charts; dropping quasi-compactness
  makes the conclusion false.
\end{lemma}

% SOURCE QUOTE PROOF: "Proof. Let \(U = \operatorname{Spec} A\) be an open
% affine subset of \(X\) on which \(f\) is regular, \(f \in A\). Then for any
% prime divisor \(Y\) which meets \(U\), \(Y \cap U\) corresponds to a minimal
% prime ideal \(\mathfrak p\) of \(A\) containing \(f\). But these are finite
% in number (II, 3.10). On the other hand, \(X - U\) is a proper closed
% subset of \(X\), so it can contain only finitely many prime divisors of
% \(X\) [\dots]"
\begin{proof}
  \uses{def:order_at_point, lem:finite_order_support_affine, lem:finite_order_support_on_affineOpen}
  Following Hartshorne II.6.1 / Stacks 02RV:
  \begin{enumerate}
  \item Pick a nonempty affine open \(U = \operatorname{Spec} R \subset X\) on
    which \(f\) is regular. Then \(R\) is a Noetherian integral domain.
  \item For each prime divisor \(Y\) meeting \(U\), the corresponding
    prime \(\mathfrak p_Y\) is a height-1 prime of \(R\); the height-1 primes
    containing \(f \cdot 1\) are precisely the minimal primes of
    \((f)\), which by Krull's Hauptidealsatz form a finite set (Stacks tag
    00KZ / Atiyah--Macdonald 11.16). Of those, \(\ord_Y(f) \ne 0\) holds
    only at the finitely many minimal primes of \((f)\) over \(R\).
  \item The remaining prime divisors \(Y\) lie in the closed complement
    \(X \setminus U\), which is a proper closed subset of \(X\); the
    Noetherian topological structure of \(X\) (Stacks 04ZE) bounds the
    irreducible components of \(X \setminus U\) of codimension one in
    \(X\) by a finite number. For each such component, the stalk
    \(\mathcal O_{X, Y}\) is a DVR (regularity in codim 1) and \(f\) has
    finite valuation; together these contribute only finitely many
    nonzero values to the order function.
  \item Combining (2) and (3) gives the finiteness of the support.
  \end{enumerate}
  \emph{Why compactness is needed.} Finite support follows because a
  compact (finite-affine-cover) scheme has only finitely many codimension-one
  prime divisors meeting the support of \(\mathrm{div}(f)\): cover \(X\) by
  finitely many affine opens \(U_1, \dots, U_N\); on each \(U_i\) step (2)
  bounds the nonzero-order prime divisors meeting \(U_i\) by the finite set of
  minimal primes of \((f)\) in the Noetherian ring \(\Gamma(U_i, \mathcal O)\);
  every prime divisor \(Y\) meets at least one \(U_i\), so the support is a
  finite union of finite sets. Quasi-compactness is genuinely necessary: a
  non-quasi-compact integral locally Noetherian scheme can carry infinitely
  many pairwise-disjoint codimension-one prime divisors outside a fixed affine
  open \(U\), so a rational function may have infinitely many zeros or poles
  and the conclusion \(\{Y \mid \ord_Y(f) \ne 0\}\) finite fails.
\end{proof}

\subsection{Affine bridge ingredients}
\label{subsec:affine_bridge_ingredients}
The affine finiteness core below rests on three Mathlib-supplied identifications
that translate scheme-level data on \(V = \operatorname{Spec} R\) into
ring-theoretic data on the Noetherian domain \(R\): the stalk of the structure
sheaf is the localization at the corresponding prime, and the function field is
the fraction field of \(R\).

\begin{lemma}[Stalk of \(\operatorname{Spec} R\) is the localization at the prime]
  \label{lem:structureSheaf_stalkIso}
  \lean{AlgebraicGeometry.StructureSheaf.stalkIso}
  \mathlibok
  \textit{Provided by Mathlib.}
  For a commutative ring \(R\) and a point \(x \in \operatorname{Spec} R\) --- i.e.\ a
  prime ideal \(\mathfrak p_x = x.\mathtt{asIdeal}\) of \(R\) --- there is an
  \(R\)-algebra isomorphism
  \[
    \mathrm{Localization.AtPrime}\,\mathfrak p_x
    \;\xrightarrow{\ \sim\ }\;
    \mathcal O_{\operatorname{Spec} R,\,x}
  \]
  between the localization of \(R\) at \(\mathfrak p_x\) and the stalk of the
  structure sheaf at \(x\).
\end{lemma}

\begin{lemma}[Function field of \(\operatorname{Spec} R\) is the fraction field]
  \label{lem:functionField_isFractionRing_of_affine}
  \lean{AlgebraicGeometry.functionField_isFractionRing_of_affine}
  \mathlibok
  \textit{Provided by Mathlib.}
  For an integral domain \(R\), the canonical map exhibits the function field of
  \(\operatorname{Spec} R\) as the fraction field of \(R\): one has
  \(\mathtt{IsFractionRing}\,R\,K(\operatorname{Spec} R)\), so
  \(K(\operatorname{Spec} R) \cong \mathrm{Frac}(R)\).
\end{lemma}

\begin{lemma}\leanok
[Affine ring-theoretic core of finite support]
  \label{lem:finite_order_support_affine}
  \lean{AlgebraicGeometry.finite_order_support_affine}
  \uses{def:order_at_point, def:prime_divisor}
  \textit{Source: Hartshorne, II.6.1 / Stacks 02RV, affine case.}
  Let \(V = \operatorname{Spec} R\) be an affine integral locally Noetherian
  scheme that is regular in codimension one, and let \(g \in K(V)^{\times}\) be a
  nonzero rational function, written \(g = a/b\) with \(a, b \in R \setminus \{0\}\).
  Then \(\{Z \in V^{(1)} \mid \ord_Z(g) \ne 0\}\) is finite.
\end{lemma}
\begin{proof}
  \uses{def:order_at_point, def:prime_divisor, def:isRegularInCodimensionOne,
    lem:structureSheaf_stalkIso, lem:functionField_isFractionRing_of_affine,
    lem:ordFrac_eq_one_of_notMem,
    lem:finite_height_one_primes_mem_or}
  Write \(R := \Gamma(V, \mathcal O_V)\), a Noetherian integral domain (\(V\) is
  integral, locally Noetherian, and affine), and identify
  \(V \cong \operatorname{Spec} R\). The scheme-to-ring bridge decomposes into
  three steps; the existing ring-theoretic count then closes the support.
  \begin{enumerate}
  \item[(i)] \emph{Prime divisor \(\leftrightarrow\) height-one prime.} A prime
    divisor \(Z\) of \(V\) (\ref{def:prime_divisor}) is a codimension-one point.
    Under \(V \cong \operatorname{Spec} R\) the point \(Z\) corresponds to a
    prime ideal \(\mathfrak p_Z \subset R\) of height one (codimension one in
    \(V\) corresponds to height one in \(R\)). Regularity in codimension one
    (\ref{def:isRegularInCodimensionOne}) makes \(R_{\mathfrak p_Z}\) a
    discrete valuation ring.
  \item[(ii)] \emph{Stalk \(\cong\) localization; function field \(\cong\)
    fraction field.} The stalk isomorphism
    \(\mathcal O_{V, Z} \cong R_{\mathfrak p_Z}\)
    (\ref{lem:structureSheaf_stalkIso}) and the function-field identification
    \(K(V) \cong \mathrm{Frac}(R)\)
    (\ref{lem:functionField_isFractionRing_of_affine}) carry the
    order-of-vanishing data from the abstract stalk and function field onto the
    explicit localization \(R_{\mathfrak p_Z} \hookrightarrow \mathrm{Frac}(R)\).
  \item[(iii)] \emph{Order \(=\) valuation.} Under these isomorphisms
    \(\ord_Z(g)\) is the \(\mathfrak p_Z\)-adic valuation
    \(v_{\mathfrak p_Z}(g)\) on \(R_{\mathfrak p_Z}\). Writing \(g = a/b\) with
    \(a, b \in R \setminus \{0\}\): if neither \(a\) nor \(b\) lies in
    \(\mathfrak p_Z\) then both are units in \(R_{\mathfrak p_Z}\), so
    \(\ord_Z(g) = 0\); contrapositively
    \[
      \ord_Z(g) \ne 0 \;\Longrightarrow\; a \in \mathfrak p_Z \ \lor\
      b \in \mathfrak p_Z .
    \]
  \end{enumerate}
  Hence \(Z \mapsto \mathfrak p_Z\) injects the support
  \(\{Z \mid \ord_Z(g) \ne 0\}\) into the set of height-one primes of \(R\)
  containing \(a\) or \(b\), which is finite by
  \ref{lem:finite_height_one_primes_mem_or} (these primes are exactly the
  minimal primes of \((a \cdot b)\)). The support is therefore finite. The
  numerator/denominator decomposition \(g = a/b\) (not merely a regular \(f\))
  is essential: zeros come from \(a\), poles from \(b\), and both are captured
  by the height-one primes containing \(a\) or \(b\).

  \emph{Residue-field independence.} The correspondence is insensitive to the
  residue field \(\kappa(Z)\), so the finiteness holds for an arbitrary affine
  integral locally-Noetherian regular-in-codimension-one scheme \(V\), with no
  algebraically-closed-base hypothesis.
\end{proof}

\subsection{Ring-theoretic core: finitely many height-one primes divide a fixed
element}
\label{subsec:ring_core_finiteness}
Step (iii)\(\to\)finiteness of \ref{lem:finite_order_support_affine} bottoms out
in a purely ring-theoretic count: in a Noetherian integral domain only finitely
many height-one primes can contain a fixed nonzero element (resp.\ one of two
fixed nonzero elements). These are the axiom-clean substrate lemmas the affine
bridge consumes.

\begin{lemma}[Order one at a prime away from both numerator and denominator]
  \label{lem:ordFrac_eq_one_of_notMem}
  \lean{Ring.ordFrac_eq_one_of_notMem}
  Let \(A\) be a localization of a Noetherian integral domain \(R\) at a prime
  \(\mathfrak p\) (\(\texttt{IsLocalization.AtPrime}\,A\,\mathfrak p\)), with common
  fraction field \(K\). If \(a \notin \mathfrak p\) and \(b \notin \mathfrak p\),
  then \(\texttt{Ring.ordFrac}\,A\,\bigl(e_K a / e_K b\bigr) = 1\), where \(e_K\)
  is the localization map \(R \to K\). Contrapositively, \(\texttt{ord}_Z(g) \neq 0
  \Rightarrow a \in \mathfrak p_Z \lor b \in \mathfrak p_Z\) for \(g = a/b\) ---
  the step-(iii) reduction the affine bridge uses to inject the order-support into
  the finite height-one-prime set of \ref{lem:finite_height_one_primes_mem_or}.
\end{lemma}
\begin{proof}
  \(a, b \notin \mathfrak p\) become units in \(A\)
  (\(\texttt{IsLocalization.map\_units}\)); pushing them through the scalar tower
  \(R \to A \to K\) makes \(e_K a / e_K b\) a unit in \(A\), so
  \(\texttt{Ring.ordFrac\_of\_isUnit}\) gives order \(1\).
\end{proof}

\begin{lemma}[Finiteness of minimal primes over an ideal in a Noetherian ring]
  \label{lem:finite_minimalPrimes_mathlib}
  \lean{Ideal.finite_minimalPrimes_of_isNoetherianRing}
  \mathlibok
  \textit{Provided by Mathlib.}
  In a Noetherian commutative ring \(R\), every ideal \(I\) has only finitely
  many minimal primes: the set \(\mathrm{minimalPrimes}(I)\) is finite. This is
  the Noetherian form of Krull's Hauptidealsatz count.
\end{lemma}

\begin{lemma}
[Finiteness of height-one primes containing a fixed nonzero element]
  \label{lem:finite_height_one_primes_mem}
  \lean{Ideal.finite_setOf_isPrime_height_one_mem}
  \uses{lem:finite_minimalPrimes_mathlib}
  Let \(R\) be a Noetherian integral domain and \(a \in R\) nonzero. Then the set
  \[
    \{\mathfrak p \subseteq R \mid \mathfrak p \text{ prime},\
      \mathrm{ht}\,\mathfrak p = 1,\ a \in \mathfrak p\}
  \]
  of height-one primes containing \(a\) is finite.
\end{lemma}
\begin{proof}
  \uses{lem:finite_minimalPrimes_mathlib}
  Every height-one prime \(\mathfrak p \ni a\) is a minimal prime over the
  principal ideal \((a)\). Indeed, there exists a minimal prime
  \(\mathfrak q \subseteq \mathfrak p\) over \((a)\), and
  \(\mathfrak q \ne 0\) since \(a \ne 0\); if
  \(\mathfrak q \subsetneq \mathfrak p\), the strict chain
  \(0 \subsetneq \mathfrak q \subsetneq \mathfrak p\) gives
  \(\mathrm{ht}\,\mathfrak p \ge 2\), contradicting \(\mathrm{ht}\,\mathfrak p = 1\).
  Hence \(\mathfrak q = \mathfrak p\), so the set is contained in
  \(\mathrm{minimalPrimes}(a)\), which is finite by
  \ref{lem:finite_minimalPrimes_mathlib}.
\end{proof}

\begin{lemma}
[Finiteness of height-one primes containing one of two fixed nonzero elements]
  \label{lem:finite_height_one_primes_mem_or}
  \lean{Ideal.finite_setOf_isPrime_height_one_mem_or}
  \uses{lem:finite_height_one_primes_mem}
  Let \(R\) be a Noetherian integral domain and \(a, b \in R\) both nonzero. Then
  the set
  \[
    \{\mathfrak p \subseteq R \mid \mathfrak p \text{ prime},\
      \mathrm{ht}\,\mathfrak p = 1,\ a \in \mathfrak p \ \lor\ b \in \mathfrak p\}
  \]
  is finite. This is the \(g = a/b\) consumable form fed to the affine support
  count (\ref{lem:finite_order_support_affine}).
\end{lemma}
\begin{proof}
  \uses{lem:finite_height_one_primes_mem}
  The set is the union of the height-one primes containing \(a\) and those
  containing \(b\); each is finite by \ref{lem:finite_height_one_primes_mem}
  (applied to \(a \ne 0\) and to \(b \ne 0\)), and a union of two finite sets is
  finite.
\end{proof}

\begin{lemma}\leanok
[Open-chart finiteness of the order support]
  \label{lem:finite_order_support_on_affineOpen}
  \lean{AlgebraicGeometry.finite_order_support_on_affineOpen}
  \uses{lem:finite_order_support_affine, lem:order_eq_order_restrict, lem:primeDivisor_ofOpen_point}
  Let \(U \subseteq X\) be a nonempty affine open of an integral locally
  Noetherian scheme \(X\) regular in codimension one, and \(f \in K(X)^{\times}\).
  Then \(\{Y \in X^{(1)} \mid Y.\mathrm{point} \in U,\ \ord_Y(f) \ne 0\}\) is finite.
\end{lemma}
\begin{proof}
  \uses{lem:finite_order_support_affine, lem:order_eq_order_restrict}
  Transport along the open immersion \(U \hookrightarrow X\): a prime divisor of
  \(X\) meeting \(U\) restricts to a prime divisor of the chart \(U\)
  (\ref{lem:primeDivisor_ofOpen_point}), order is preserved
  (\ref{lem:order_eq_order_restrict}), and the pulled-back function on the chart
  has finite nonzero-order support by \ref{lem:finite_order_support_affine}.
  The round-trip \(\mathrm{ofOpen} \circ \mathrm{restrictToOpen} = \mathrm{id}\)
  identifies the two index sets.
\end{proof}

\begin{definition}
\leanok
  \uses{def:codim1_cycles, def:order_at_point, lem:rationalMap_order_finite_support}
[Principal divisor of a rational function]
  \label{def:principal_divisor}
  
  \lean{AlgebraicGeometry.Scheme.WeilDivisor.principal}
  % SOURCE: [Hartshorne], II.6, Lemma~6.1 + the definition of \(\operatorname{div}(f) = (f)\),
  % p.~131 (read from references/hartshorne-algebraic-geometry.pdf, PDF page 148).
  % SOURCE QUOTE: "Lemma 6.1. Let \(X\) satisfy \((\ast)\), and let \(f \in K^\ast\)
  % be a nonzero function on \(X\). Then \(v_Y(f) = 0\) for all except finitely
  % many prime divisors \(Y\)."
  % SOURCE QUOTE: "Definition. Let \(X\) satisfy \((\ast)\) and let \(f \in K^\ast\).
  % We define the \emph{divisor of \(f\)}, denoted \((f)\), by
  % \((f) = \sum v_Y(f) \cdot Y\), where the sum is taken over all prime
  % divisors \(Y\) of \(X\). By the lemma, this is a finite sum, hence is a
  % divisor. Any divisor which is equal to the divisor of a function is
  % called a \emph{principal divisor}."
  \textit{Source: Hartshorne, II.6, Lemma~6.1 and the following definition,
  p.~131.}
  Let \(X\) satisfy \((*)\). For a nonzero rational function \(f \in K(X)^{\times}\),
  Hartshorne's Lemma~6.1 asserts that \(\ord_Y(f) = 0\) for all but finitely
  many prime divisors \(Y\) on \(X\); hence the formal sum
  \[
    \mathrm{div}(f) \;:=\; \sum_{\substack{Y \subset X \\ \text{prime divisor}}}
    \ord_Y(f) \cdot Y \;\in\; \mathrm{Div}(X)
  \]
  has finite support and is a well-defined Weil divisor, called the
  \emph{principal divisor of \(f\)}. On a smooth proper curve \(C\) over \(\bar k\)
  the same definition specialises to
  \(\mathrm{div}(f) = \sum_{P \in C \text{ closed}} \ord_P(f) \cdot [P]\), the formal
  sum over closed points.
\end{definition}
\begin{lemma}

\leanok
[Coefficient of a principal divisor is the order]
  \label{lem:principal_apply}
  \lean{AlgebraicGeometry.Scheme.WeilDivisor.principal_apply}
  \uses{def:principal_divisor, def:order_at_point}
  Let \(X\) satisfy \((*)\), \(f \in K(X)^{\times}\), and \(Y\) a prime divisor
  of \(X\). Then the coefficient of the principal divisor \(\mathrm{div}(f)\)
  (\ref{def:principal_divisor}) at \(Y\) is the order \(\ord_Y(f)\)
  (\ref{def:order_at_point}). This is the basic unfolding of the
  finite-support packaging used to define \(\mathrm{div}(f)\).
\end{lemma}
\begin{proof}
  Proved directly in Lean (from \(\texttt{Finsupp.ofSupportFinite\_coe}\)).
\end{proof}
\begin{lemma}

\leanok
[Principal divisor of the constant function one]
  \label{lem:principal_one}
  \lean{AlgebraicGeometry.Scheme.WeilDivisor.principal_one}
  \uses{def:principal_divisor, lem:principal_apply, lem:order_one}
  Let \(X\) satisfy \((*)\). The principal divisor of the constant function
  \(1 \in K(X)^{\times}\) is the zero divisor: \(\mathrm{div}(1) = 0\).
\end{lemma}
\begin{proof}
  Proved directly in Lean: by \ref{lem:principal_apply} each coefficient is
  \(\ord_Y(1)\), which is \(0\) by \ref{lem:order_one}.
\end{proof}
\begin{lemma}

\leanok
[Principal divisor of the inverse is the negation]
  \label{lem:principal_inv}
  \lean{AlgebraicGeometry.Scheme.WeilDivisor.principal_inv}
  \uses{def:principal_divisor, lem:principal_apply, lem:order_inv}
  Let \(X\) satisfy \((*)\) and let \(f \in K(X)^{\times}\). Then the principal
  divisor of \(f^{-1}\) is the negation of the principal divisor of \(f\):
  \[
    -\,\mathrm{div}(f) \;=\; \mathrm{div}(f^{-1}) \;\in\; \mathrm{Div}(X).
  \]
  Equivalently, the zeros of \(f^{-1}\) are exactly the poles of \(f\) and vice
  versa. This is the divisor-level form of the per-prime-divisor valuation
  identity \(v_Y(f^{-1}) = -v_Y(f)\), and it is the substrate for the
  homomorphism law (\ref{thm:principal_hom}), the linear-equivalence symmetry,
  and the zeros--poles bookkeeping in \ref{thm:principal_deg_zero}.
\end{lemma}
\begin{proof}
  Proved directly in Lean coefficientwise: by \ref{lem:principal_apply} the
  coefficient of each side at a prime divisor \(Y\) is \(-\ord_Y(f)\),
  respectively \(\ord_Y(f^{-1})\), and these agree by the order-of-inverse
  identity \(\ord_Y(f^{-1}) = -\ord_Y(f)\) (\ref{lem:order_inv}).
\end{proof}
\begin{theorem}
\leanok
[The principal-divisor map is a group homomorphism]
  \label{thm:principal_hom}
  \lean{AlgebraicGeometry.Scheme.WeilDivisor.principal_hom}
  \uses{def:principal_divisor, lem:order_mul_of_ne_zero, lem:order_units_inv,
    lem:order_one}
  % SOURCE: [Hartshorne], II.6, the paragraph after the principal-divisor
  % definition, p.~131 (read from
  % references/hartshorne-algebraic-geometry.pdf, PDF page 148).
  % SOURCE QUOTE: "Note that if \(f, g \in K^\ast\), then \((f/g) = (f) - (g)\)
  % because of the properties of valuations. Therefore sending a function
  % \(f\) to its divisor \((f)\) gives a homomorphism of the multiplicative
  % group \(K^\ast\) to the additive group \(\operatorname{Div} X\), and the
  % image, which consists of the principal divisors, is a subgroup of
  % \(\operatorname{Div} X\)."
  \textit{Source: Hartshorne, II.6, p.~131 (group-homomorphism remark).}
  Let \(X\) satisfy \((*)\). The principal-divisor map
  \(\mathrm{div} \colon K(X)^{\times} \to \mathrm{Div}(X)\) is a group homomorphism from the
  multiplicative group of nonzero rational functions to the additive group
  of Weil divisors. Concretely
  \[
    \mathrm{div}(fg) \;=\; \mathrm{div}(f) + \mathrm{div}(g),
    \qquad
    \mathrm{div}(f^{-1}) \;=\; -\mathrm{div}(f),
    \qquad
    \mathrm{div}(1) \;=\; 0.
  \]
\end{theorem}
\begin{proof}
  Each prime divisor \(Y\) on \(X\) contributes a discrete valuation \(v_Y\) on
  \(K(X)\) (\ref{def:principal_divisor}); by the standard DVR axioms,
  \(v_Y(fg) = v_Y(f) + v_Y(g)\) and \(v_Y(f^{-1}) = -v_Y(f)\) for $f, g \in
  K(X)^{\times}$, and $v_Y(1) = 0$. Summing these identities over all prime
  divisors of \(X\) (a sum that has finite support on each factor by
  \ref{lem:rationalMap_order_finite_support}, so the support of the sum is the union
  of the two finite supports and remains finite) gives the three displayed
  identities.
\end{proof}
\begin{theorem}
\leanok
[Principal divisors have degree zero on a complete nonsingular curve]
  \label{thm:principal_deg_zero}
  \lean{AlgebraicGeometry.Scheme.WeilDivisor.principal_degree_zero}
  \uses{def:principal_divisor, def:divisor_degree}
  % SOURCE: [Hartshorne], II.6, Corollary~6.10, p.~138 (read from
  % references/hartshorne-algebraic-geometry.pdf, PDF page 155).
  % Stacks pointer: tag 0BE3.
  % SOURCE QUOTE: "Corollary 6.10. \emph{A principal divisor on a complete
  % nonsingular curve \(X\) has degree zero. Consequently the degree function
  % induces a surjective homomorphism $\deg \colon \operatorname{Cl} X \to
  % \mathbf Z$.}"
  \textit{Source: Hartshorne, II.6, Corollary~6.10, p.~138.}
  Let \(C\) be a smooth proper curve over an algebraically closed field
  \(\bar k\). For every nonzero rational function \(f \in K(C)^{\times}\) the
  associated principal divisor has degree zero:
  \[
    \deg(\mathrm{div}(f)) \;=\; 0 \;\in\; \mathbb Z.
  \]
\end{theorem}
% SOURCE QUOTE PROOF: "PROOF. Let \(f \in K(X)^\ast\). If \(f \in k\), then
% \((f) = 0\), so there is nothing to prove. If \(f \notin k\), then the
% inclusion of fields \(k(f) \subset K(X)\) induces a finite morphism
% \(\varphi \colon X \to \mathbf P^1\). It is a morphism by (I, 6.12). We
% have \((f) = \varphi^\ast(\{0\} - \{\infty\})\). Since \(\{0\} - \{\infty\}\)
% is a divisor of degree 0 on \(\mathbf P^1\), we conclude that \((f)\) has
% degree 0 on \(X\)."
\begin{proof}
  Following Hartshorne's Corollary~6.10. If \(f \in \bar k \subset K(C)\) is
  a constant scalar, then \(f\) has no zeros or poles, so \(\mathrm{div}(f) = 0\) and
  \(\deg(\mathrm{div}(f)) = 0\). Otherwise the inclusion of fields
  \(\bar k(f) \subset K(C)\) exhibits \(K(C)\) as a finite extension of the
  rational function field \(\bar k(f) \cong \bar k(t)\), so the corresponding
  morphism \(\varphi \colon C \to \mathbb P^1_{\bar k}\) is finite (by the
  smooth-proper-curve / function-field correspondence on which a curve is
  determined by its function field, classical I.6.12). The principal
  divisor \(\mathrm{div}(f)\) equals the pullback \(\varphi^{*}([0] - [\infty])\)
  along \(\varphi\) of the degree-\(0\) divisor \([0] - [\infty]\) on
  \(\mathbb P^1_{\bar k}\); since pullback along a finite morphism
  multiplies degrees by \(\deg(\varphi)\), we conclude
  $\deg(\mathrm{div}(f)) = \deg(\varphi) \cdot \deg([0] - [\infty]) = \deg(\varphi)
  \cdot 0 = 0$.
  The two auxiliary statements ``the inclusion \(\bar k(f) \subset K(C)\)
  defines a finite morphism \(C \to \mathbb P^1\)'' and ``pullback along
  a finite morphism of curves multiplies degrees by the degree of the
  morphism'' (Hartshorne's Proposition~II.6.9) are sub-lemmas of this
  chapter, and are depended on by the \texttt{RR.3}/\texttt{RR.4} chain.
\end{proof}
\section{Positive part of a Weil divisor}
\begin{definition}


\leanok
[Positive part of a Weil divisor]
  \label{def:WeilDivisor_positivePart}
  \lean{AlgebraicGeometry.Scheme.WeilDivisor.positivePart}
  \uses{def:codim1_cycles}
  \textit{Source: Project-bespoke (Hartshorne II.6.10 phrasing).}
  Let \(X\) be a scheme satisfying the codim-1 hypotheses of
  \ref{def:codim1_cycles}. For a Weil divisor
  \(D = \sum_Y n_Y \cdot Y \in \mathrm{Div}(X)\), the \emph{positive part}
  \((D)_0 \in \mathrm{Div}(X)\) is the Weil divisor
  \[
    (D)_0 \;:=\; \sum_{Y \,:\, n_Y > 0} n_Y \cdot Y.
  \]
  Concretely, \((D)_0\) is obtained from \(D\) by replacing each coefficient
  \(n_Y\) with \(\max(n_Y, 0)\). Equivalently, in the \(\Finsupp{\mathrm{PrimeDivisor}(X)}{\mathbb Z}\)
  presentation, the positive part is obtained by the pointwise lattice
  join with the zero divisor — \(\,(D)_0 = D \sqcup 0\,\) — equivalently
  by mapping coefficients through \(n \mapsto n \sqcup 0\):
  \(\,(D)_0 = \mathtt{Finsupp.mapRange}\;(n \mapsto n \sqcup 0)\;D\). The
  complementary \emph{negative part}
  \((D)_\infty\) is defined by \((D)_\infty := (-D)_0\), so that
  \(D = (D)_0 - (D)_\infty\) is the canonical decomposition into an
  effective (non-negative-coefficient) divisor minus an effective divisor.
  On a smooth proper curve \(C\) over \(\bar k\), where the prime divisors
  are the closed points and the coefficients are integers, the positive part
  reduces to keeping only the closed points whose multiplicity is positive
  — i.e. precisely the zeros of any rational function whose principal
  divisor is \(D\), in the case \(D = \mathrm{div}(f)\).
\end{definition}
\begin{lemma}

\leanok
[Positive part of the zero divisor]
  \label{lem:positivePart_zero}
  \lean{AlgebraicGeometry.Scheme.WeilDivisor.positivePart_zero}
  \uses{def:WeilDivisor_positivePart}
  The positive part of the zero divisor is zero: \((0)_0 = 0\).
\end{lemma}
\begin{proof}
  Proved directly in Lean (from \(\texttt{Finsupp.mapRange\_zero}\)).
\end{proof}
\begin{lemma}

\leanok
[Positive part of a one-point divisor]
  \label{lem:positivePart_single}
  \lean{AlgebraicGeometry.Scheme.WeilDivisor.positivePart_single}
  \uses{def:WeilDivisor_positivePart, def:prime_divisor}
  For a prime divisor \(Y\) and \(n \in \mathbb Z\), the positive part of the
  one-point-supported divisor \(n \cdot Y\) is \(\max(n, 0) \cdot Y\):
  \((n \cdot Y)_0 = \max(n, 0) \cdot Y\).
\end{lemma}
\begin{proof}
  Proved directly in Lean (from \(\texttt{Finsupp.mapRange\_single}\), with
  \(\max(0, 0) = 0\) the zero-preservation witness).
\end{proof}
\begin{lemma}

\leanok
[Degree of the positive part as a sum of capped coefficients]
  \label{lem:degree_positivePart_eq_sum_max}
  \lean{AlgebraicGeometry.Scheme.WeilDivisor.degree_positivePart_eq_sum_max}
  \uses{def:WeilDivisor_positivePart, def:divisor_degree}
  For a Weil divisor \(D \in \mathrm{Div}(X)\), the degree of the positive part
  equals the sum of the capped coefficients:
  \(\deg\bigl((D)_0\bigr) = \sum_Y \max(n_Y, 0)\), where \(D = \sum_Y n_Y \cdot Y\).
\end{lemma}
\begin{proof}
  Proved directly in Lean (unfold \ref{def:WeilDivisor_positivePart} and
  \ref{def:divisor_degree}; one step via \(\texttt{Finsupp.sum\_mapRange\_index}\)).
\end{proof}
\begin{lemma}

\leanok
[Degree splits into positive and negative parts]
  \label{lem:degree_eq_degree_positivePart_sub}
  \lean{AlgebraicGeometry.Scheme.WeilDivisor.degree_eq_degree_positivePart_sub}
  \uses{def:WeilDivisor_positivePart, def:divisor_degree,
    lem:degree_positivePart_eq_sum_max}
  For every Weil divisor \(D \in \mathrm{Div}(X)\), the degree decomposes as the
  degree of its positive part minus the degree of the positive part of its
  negation:
  \[
    \deg(D) \;=\; \deg\bigl((D)_0\bigr) \;-\; \deg\bigl((-D)_0\bigr).
  \]
  Writing \((D)_\infty := (-D)_0\) for the negative part
  (\ref{def:WeilDivisor_positivePart}), this is the degree-level form of the
  canonical decomposition \(D = (D)_0 - (D)_\infty\) of a divisor into a
  difference of two effective divisors --- the divisor of zeros minus the
  divisor of poles in the case \(D = \mathrm{div}(f)\).
\end{lemma}
\begin{proof}
  Proved directly in Lean. By \ref{lem:degree_positivePart_eq_sum_max} both
  positive-part degrees are sums of capped coefficients; the
  \((-D)\)-sum is turned into a \(D\)-sum via \(\texttt{Finsupp.sum\_neg\_index}\),
  and the two sums are combined into one, which is then closed coefficientwise
  by the integer identity \(n = \max(n,0) - \max(-n,0)\).
\end{proof}
\begin{lemma}
[Capped sum equals the sum over positive-coefficient points]
  \label{lem:finsupp_sum_max_zero_eq_sum_filter_pos}
  \lean{Finsupp.sum_max_zero_eq_sum_filter_pos}
  \uses{lem:degree_positivePart_eq_sum_max}
  For a finitely supported function \(D \colon \alpha \to_0 \mathbb Z\),
  \(\sum_a \max(D_a, 0) = \sum_{a \,:\, 0 < D_a} D_a\), the sum running over
  the points of positive coefficient. Negative-coefficient points contribute
  \(0\) to the cap and drop out; positive-coefficient points contribute
  \(D_a\) to both sides. This is the combinatorial identity isolating the
  effective (zeros) part of a divisor in the degree computation of
  \ref{lem:degree_positivePart_eq_sum_max}.
\end{lemma}
\begin{proof}
  Proved directly in Lean (a \(\texttt{Finset.sum\_filter}\) rewrite followed
  by a pointwise case-split on the sign of \(D_a\)).
\end{proof}
\begin{lemma}

\leanok
[A single simple zero forces positive-part degree at least one]
  \label{lem:one_le_degree_positivePart_principal_of_order_one}
  \lean{AlgebraicGeometry.Scheme.WeilDivisor.one_le_degree_positivePart_principal_of_order_one}
  \uses{lem:degree_positivePart_eq_sum_max, lem:principal_apply,
    def:principal_divisor, def:order_at_point}
  Let \(X\) satisfy \((*)\) and let \(f \in K(X)^{\times}\). If \(f\) has order
  \(1\) at some prime divisor \(Y_0\) (\(\ord_{Y_0}(f) = 1\)), then the degree
  of the positive part of the principal divisor \(\mathrm{div}(f)\) is at least
  one: \(\deg\bigl((\mathrm{div}(f))_0\bigr) \geq 1\).
\end{lemma}
\begin{proof}
  Proved directly in Lean: rewrite the degree of the positive part as a sum of
  capped coefficients (\ref{lem:degree_positivePart_eq_sum_max}); isolate the
  \(Y_0\) term, which by \ref{lem:principal_apply} equals \(\max(1, 0) = 1\),
  and bound the remaining sum below by \(0\) since each capped coefficient is
  non-negative.
\end{proof}
\begin{definition}

\leanok
[Projective line over \(\bar k\) is locally Noetherian]
  \label{def:projectiveLineBar_isLocallyNoetherian}
  \lean{AlgebraicGeometry.Scheme.WeilDivisor.instIsLocallyNoetherianProjectiveLineBar}
  \uses{def:codim1_cycles}
  The scheme \(\mathbb P^1_{\bar k}\) underlying \(\texttt{ProjectiveLineBar}\,\bar k\)
  is locally Noetherian. This is the typeclass needed to run the order /
  principal / degree machinery on \(\mathbb P^1_{\bar k}\).
\end{definition}
\begin{proof}
  Proved directly in Lean: the structure morphism \(\mathbb P^1_{\bar k} \to
  \operatorname{Spec}\bar k\) is proper, hence locally of finite type; the base
  \(\operatorname{Spec}\bar k\) is locally Noetherian (a field is Noetherian);
  a locally-finite-type morphism over a locally Noetherian base has a locally
  Noetherian source.
\end{proof}
\begin{theorem}

\leanok
[Projective line over \(\bar k\) is regular in codimension one]
  \label{thm:projectiveLineBar_isRegularInCodimOne}
  \lean{AlgebraicGeometry.Scheme.WeilDivisor.isRegularInCodimOneProjectiveLineBar}
  \uses{def:isRegularInCodimensionOne}
  The scheme \(\mathbb P^1_{\bar k}\) underlying \(\texttt{ProjectiveLineBar}\,\bar k\)
  satisfies \(\texttt{IsRegularInCodimensionOne}\)
  (\ref{def:isRegularInCodimensionOne}): every codimension-one stalk is a
  discrete valuation ring.
\end{theorem}
\begin{proof}
  By the standard ``smooth of relative dimension one over \(\bar k\)
  \(\Rightarrow\) regular \(\Rightarrow\) discrete valuation ring at every
  codimension-one point'' chain, reduced through the two-chart affine cover of
  \(\mathbb P^1_{\bar k}\), each chart being \(\operatorname{Spec}\) of a
  polynomial ring (a principal ideal domain, hence Dedekind), so localising at
  a height-one prime yields a discrete valuation ring.
\end{proof}
\begin{lemma}

\leanok
[Degree of the positive part of \(\mathrm{div}(f)\) for a non-constant rational
function equals the function-field extension degree]
  \label{lem:degree_positivePart_principal_eq_finrank}
  \lean{AlgebraicGeometry.Scheme.WeilDivisor.degree_positivePart_principal_eq_finrank}
  \uses{def:WeilDivisor_positivePart, def:principal_divisor,
        def:order_at_point, def:divisor_degree, def:prime_divisor,
        lem:degree_positivePart_eq_sum_max,
        lem:finsupp_sum_max_zero_eq_sum_filter_pos,
        lem:one_le_degree_positivePart_principal_of_order_one,
        lem:principal_apply, lem:order_pow_of_ne_zero,
        lem:degree_single, lem:positivePart_single,
        def:projectiveLineBar_isLocallyNoetherian,
        thm:projectiveLineBar_isRegularInCodimOne}
  % SOURCE: [Hartshorne], II.6, Proposition~6.9, p.~137 (multiplicativity
  % of degree under finite pullback) (read from
  % references/hartshorne-algebraic-geometry.pdf, PDF page 154).
  % Stacks pointer: tag 02RW (pullback of divisors under finite
  % morphisms), tag 0BE6 (degree under finite morphism).
  % SOURCE QUOTE: "Proposition 6.9. \emph{Let \(f \colon X \to Y\) be a finite
  % morphism of nonsingular curves. Then for any divisor \(D\) on \(Y\),
  % \(\deg(f^\ast D) = [K(X) : K(Y)] \cdot \deg D\).}"
  \textit{Source: Hartshorne, II.6, Proposition~6.9, p.~137.}
  Let \(C\) be a smooth proper geometrically irreducible curve over an
  algebraically closed field \(\bar k\), and let
  \(\varphi \colon C \to \mathbb P^1_{\bar k}\) be a non-constant morphism.
  Let \(t_\infty \in K(\mathbb P^1)\) be a local parameter at \(\infty \in \mathbb P^1\)
  (equivalently \(t_\infty = 1/u\) for the standard affine coordinate \(u\) on
  \(\mathbb P^1 \setminus \{\infty\}\)). Write \(f := \varphi^{\#} t_\infty
  \in K(C)\) for the pullback of \(t_\infty\) along \(\varphi\). Then \(f\) is
  nonzero (because \(\varphi\) is non-constant; equivalently \(\varphi^{-1}(\infty)\)
  is a proper closed subscheme of \(C\)) and the degree of the positive
  part of the principal divisor of \(f\) equals the function-field
  extension degree:
  \[
    \deg\bigl((\mathrm{div}(f))_0\bigr)
    \;=\; \bigl[\,K(C) \,:\, \bar k(\mathbb P^1)\,\bigr]
    \;=\; \mathrm{Module.finrank}_{\bar k(\mathbb P^1)} K(C).
  \]
  Equivalently, \((\mathrm{div}(f))_0 = \varphi^{*}[\infty]\) (the effective
  preimage divisor of \(\infty\)) and \(\deg(\varphi^{*}[\infty]) = \deg(\varphi)\).
\end{lemma}
% SOURCE QUOTE PROOF: "PROOF. The morphism \(f \colon X \to Y\) is finite,
% so locally on \(Y\) corresponds to a finite extension of Dedekind domains
% \(A \subset B\). For each closed point \(P \in Y\) the ramification--inertia
% formula gives \(\sum_{Q | P} e_Q f_Q = [K(X) : K(Y)]\), where the sum runs
% over closed points \(Q \in X\) above \(P\). Summing this identity over
% \(P \in Y\) using the affine cover and noting that \(D = \sum_P n_P [P]\)
% gives \(\deg(f^\ast D) = \sum_{P} n_P \sum_{Q | P} e_Q f_Q
% = \sum_P n_P [K(X) : K(Y)] = [K(X) : K(Y)] \deg(D)\)."
\begin{proof}
  Following Hartshorne's Proposition~II.6.9 specialised to the divisor
  \(D = [\infty]\) on \(\mathbb P^1_{\bar k}\). A non-constant morphism
  between smooth proper curves over \(\bar k\) is automatically finite (the
  image is a closed irreducible subset that is not a point and not empty,
  hence all of \(\mathbb P^1\); a surjective morphism between smooth proper
  irreducible curves of equal dimension is proper + quasi-finite, hence
  finite; cf.\ Hartshorne II.6.8). Hence the function-field extension
  \(\bar k(\mathbb P^1) \hookrightarrow K(C)\) is finite of degree
  \(d := [K(C) : \bar k(\mathbb P^1)] = \dim_{\bar k(\mathbb P^1)} K(C)\).
  Pick an affine open \(\mathrm{Spec}(A) \subset \mathbb P^1_{\bar k}\) containing
  \(\infty\); explicitly take \(A = \bar k[t_\infty]\) the standard chart at
  \(\infty\), with \(\infty\) corresponding to the maximal ideal
  \(\mathfrak m_\infty = (t_\infty) \subset A\). The preimage
  \(\mathrm{Spec}(B) := \varphi^{-1}(\mathrm{Spec}(A))\) is an affine open of \(C\)
  with \(B\) a finite (and integrally closed, by \(C\) smooth) extension of
  \(A\); both are Dedekind domains. Apply
  \texttt{Ideal.sum\_ramification\_inertia} (Mathlib,
  \texttt{Mathlib.NumberTheory.RamificationInertia.Basic}) to the
  ring-theoretic extension \(A \to B\) at the maximal ideal
  \(\mathfrak m_\infty\):
  \[
    \sum_{Q \,\colon\, \mathfrak m_\infty \subset Q} e_{Q | \mathfrak m_\infty} \cdot f_{Q | \mathfrak m_\infty}
    \;=\; \dim_{A / \mathfrak m_\infty}\!\bigl(B / \mathfrak m_\infty B\bigr)
    \;=\; \dim_{\bar k(\mathbb P^1)} K(C) \;=\; d,
  \]
  where the second equality uses
  \texttt{Ideal.finrank\_quotient\_map} (which identifies the residue-field
  quotient with the function-field extension degree for a Dedekind extension
  at a maximal ideal). The closed points \(Q \subset \mathrm{Spec}(B)\) above
  \(\mathfrak m_\infty\) are precisely the closed points
  \(P \in \varphi^{-1}(\infty) \subset C\); the ramification index
  \(e_{Q | \mathfrak m_\infty}\) is the order \(\ord_P(\varphi^{\#} t_\infty)\)
  of the pulled-back local parameter at \(P\); the residue degree
  \(f_{Q | \mathfrak m_\infty}\) is \([\kappa(P) : \bar k] = 1\) (every closed
  point of \(C\) over \(\bar k\) has residue field \(\bar k\)). Therefore
  \[
    \sum_{P \,\in\, \varphi^{-1}(\infty)} \ord_P(\varphi^{\#} t_\infty) \;=\; d
    \;=\; \mathrm{Module.finrank}_{\bar k(\mathbb P^1)} K(C).
  \]
  The left-hand side equals the degree of the positive part of
  \(\mathrm{div}(\varphi^{\#} t_\infty)\), because: at every closed point
  \(P \in \varphi^{-1}(\infty)\) the order is \(\geq 1\) (pullback of a local
  parameter under a finite morphism is itself a uniformiser-power, hence
  has positive order at every preimage); at every closed point
  \(P \in C \setminus \varphi^{-1}(\infty)\) the rational function
  \(\varphi^{\#} t_\infty\) is regular and nonzero (because \(t_\infty\) is
  regular and nonzero on \(\mathbb P^1 \setminus \{\infty\}\)), hence the
  order at such a \(P\) is \(\geq 0\); the points where the order is
  \emph{strictly} positive are exactly the points of \(\varphi^{-1}(\infty)\).
  Therefore
  \(\,(\mathrm{div}(\varphi^{\#} t_\infty))_0 = \sum_{P \in \varphi^{-1}(\infty)}
  \ord_P(\varphi^{\#} t_\infty) \cdot [P] = \varphi^{*}[\infty]\,\), and
  taking degree completes the proof.
\end{proof}
\section{Divisor class group}
We follow Hartshorne II.6 (pp.~131--132) and define the divisor class group
of a curve as the cokernel of the principal-divisor homomorphism. This is
the \emph{Picard group} of \(C\) in the curve-specific sense; the line-bundle
interpretation \(\Cl(C) \cong \Pic(C)\) requires the construction of
\(\mathcal O_C(D)\), which lives in the sibling chapter \texttt{RR.3} and is
out of scope here.
\begin{definition}
\leanok
[Linear equivalence of divisors]
  \label{def:linear_equivalence}
  \lean{AlgebraicGeometry.Scheme.WeilDivisor.LinearEquivalence}
  \uses{def:principal_divisor, thm:principal_hom, def:codim1_cycles}
  % SOURCE: [Hartshorne], II.6, p.~131 (definition of linear equivalence
  % and of \(\operatorname{Cl} X\)) (read from
  % references/hartshorne-algebraic-geometry.pdf, PDF page 148).
  % SOURCE QUOTE: "Definition. Let \(X\) satisfy \((\ast)\). Two divisors \(D\)
  % and \(D'\) are said to be \emph{linearly equivalent}, written \(D \sim D'\),
  % if \(D - D'\) is a principal divisor. The group \(\operatorname{Div} X\) of
  % all divisors divided by the subgroup of principal divisors is called
  % the \emph{divisor class group} of \(X\), and is denoted by
  % \(\operatorname{Cl} X\)."
  \textit{Source: Hartshorne, II.6, p.~131 (linear equivalence and
  \(\operatorname{Cl} X\)).}
  Let \(X\) satisfy \((*)\). Two Weil divisors \(D, D' \in \mathrm{Div}(X)\) are
  \emph{linearly equivalent}, written \(D \sim D'\), if and only if there
  exists \(f \in K(X)^{\times}\) with
  \[
    D - D' \;=\; \mathrm{div}(f) \;\in\; \mathrm{Div}(X).
  \]
  The relation \(\sim\) is an equivalence relation on \(\mathrm{Div}(X)\) (the
  reflexivity, symmetry, and transitivity all transport from the analogous
  identities \(\mathrm{div}(1) = 0\), \(\mathrm{div}(f^{-1}) = -\mathrm{div}(f)\), $\mathrm{div}(fg) =
  \mathrm{div}(f) + \mathrm{div}(g)$ via \ref{thm:principal_hom}); the corresponding
  quotient group \(\Cl(X) := \mathrm{Div}(X) / \mathrm{im}(\div)\) is the
  \emph{divisor class group} of \(X\). On a smooth proper curve \(C\) over
  \(\bar k\), \ref{thm:principal_deg_zero} implies that the degree map
  \(\deg \colon \mathrm{Div}(C) \to \mathbb Z\) descends to a homomorphism
  \(\deg \colon \Cl(C) \to \mathbb Z\).
\end{definition}
\section{Out of scope}
The following items are deliberately \emph{not} part of \texttt{RR.1} and
are sequenced into the sibling sub-build chapters:
\begin{itemize}
  \item The line-bundle construction \(\mathcal O_C(D)\) associated to a Weil
    divisor \(D\), the calculation of its global sections, and the
    identification \(\mathcal O_C(D) \cong \mathcal O_C(D')\) when
    \(D \sim D'\), lives in \texttt{RR.3} (sibling chapter
    \texttt{RiemannRoch\_OcOfD.tex}, to be added).
  \item The Riemann--Roch dimension formula
    \(\ell(D) - \ell(K - D) = \deg(D) + 1 - g\) and its genus-\(0\) specialisation
    \(\ell(D) = \deg(D) + 1\) for \(\deg(D) \geq 0\) live in \texttt{RR.2}
    (sibling chapter \texttt{RiemannRoch\_RRFormula.tex}, to be added).
  \item The classification ``a smooth proper geometrically irreducible
    curve over \(\bar k\) with \(g = 0\) is isomorphic to \(\mathbb P^1\)''
    (Hartshorne Example~IV.1.3.5) is the headline of \texttt{RR.4}
    (sibling chapter \texttt{RiemannRoch\_RationalIsoP1.tex}, to be added)
    and assembles \texttt{RR.1}/\texttt{RR.2}/\texttt{RR.3} into the
    \texttt{Proj.fromOfGlobalSections} morphism into \(\mathbb P^1\) together
    with its degree-\(1\) identification.
  \item \emph{Cartier divisors} (locally principal divisors on a possibly
    singular scheme) are out of scope: the project's curves are smooth, so
    Weil and Cartier divisors coincide, and the simpler Weil form suffices.
    A future Mathlib upstream PR could introduce the Cartier formalism.
  \item \emph{Higher-codimension cycles} (the Chow group
    \(\bigoplus_i \mathrm{CH}^i(X)\)) are out of scope: this chapter is
    strictly curve-focused.
  \item The proof of \ref{thm:principal_deg_zero} factors through two
    auxiliary statements (``finite morphism induced by a non-constant
    rational function'' and ``multiplicativity of degree under finite
    pullback'', Hartshorne II.6.9) which are themselves separable
    sub-lemmas of \texttt{RR.1}; they are flagged in the proof body but
    deliberately not unfolded as standalone declarations here.
\end{itemize}
